\section{Research priorities and source scope}
\label{sec:outlook}

For an initial random-$k$-SAT project, three next steps are concrete.
First, review the supplied threshold proof in detail, checking its stopping-time
and model-transfer arguments and the dependencies of the convergence criterion.
The newer window and $4.268$ upper-bound manuscripts should be assessed
separately, including any claimed formal artifacts. Second, run the modest
CPU feasibility pilot in Section~\ref{sec:scale}, keeping search and complete
decision results separate. Third, if the geometric analogy is pursued,
specify a relaxation or discrete invariant that survives the obstruction in
Proposition~\ref{prop:lp} before attempting a concentration theorem.

The algorithmic frontier is best expressed as a collection of questions.
How far into the satisfiable phase can one guarantee polynomial-time witness
search? Which computational models explain the observed barrier? Can random
UNSAT formulas at constant density be efficiently certified by stronger
methods? For small fixed $k$, how precisely can the threshold location and
finite-size bias be determined? The reviewed sources address parts of these
questions; the recent preprints do not supply a complete practical solution.

\subsection{The supplied inventory and added literature}
The repository originally supplied four PDFs. \texttt{TR26-229.pdf} is
Carenini's October 2026 random-$k$-SAT preprint. \texttt{iau005.pdf} is
Amelunxen--Lotz--McCoy--Tropp, and \texttt{sharp-phase.pdf} is Lotz--Tropp.
The file \texttt{montanari-phase.pdf} is actually Montanari and Wang's
\emph{Phase Transitions for Feature Learning in Neural Networks}
\citep{montanari2026}; it is not a second random-$k$-SAT background paper.
It has therefore not been used as evidence for a SAT theorem. The background
coverage here is supplied by the additional primary sources in the bibliography.

Literature and software results were checked through 10 October 2026.
The exposition revision of 11 October adds the elementary clause-count
proof and its teaching-note reference, and expands the clause-process definition.
The coverage is targeted, not a claim of an exhaustive literature search.
The source registry records stable S identifiers, versions, local originals,
precise theorem/table locations and the scope of each reading. Checking a
stated theorem and its hypotheses is distinguished from checking its proof;
reported computational and Lean artifacts were not rerun.

The article, bibliography and P001--P002 fragments are editable LaTeX. The repository
command \texttt{make report} produces the PDF. The project status and source
manifest identify the dated version; earlier template material is archived.
No new solver measurements are included. Consequently the immediate missing
evidence for a hardware-specific size recommendation is the local pilot,
not another extrapolation from an unrelated benchmark.
